



\documentclass[11pt,a4paper]{article}
\usepackage[margin=2.5cm]{geometry}
\usepackage[utf8]{inputenc}
\usepackage[T1]{fontenc}
\usepackage[french]{babel}
\usepackage{enumitem}
\usepackage{hyperref}


\begin{document}
	\pagestyle{empty}
	
	\noindent
	Friedly WOLI \\  % \hfill Airbus \\
	Toulouse, 31100  \\
	0767339290 \\
	friedlysedjro@gmail.com 
	
	\begin{flushright}
		
	
	\end{flushright}
	
	
	\bigskip
	
	
	
	
	
	
	\textbf{Objet :} Candidature spontanée – Alternance en simulation mécanique \\ \\
	Madame, Monsieur\\
	

			
			Actuellement étudiant en troisième année de Licence Mécanique Énergétique à l’Université Paul Sabatier, je suis à la recherche d’une alternance en modélisation et simulation mécanique afin de contribuer à des projets innovants dans le domaine de l’ingénierie mécanique.\\
			
			Ma formation m’a permis d’acquérir de solides bases en mécanique des solides, résistance des matériaux, mécanique des fluides et modélisation mathématique des systèmes physiques. J’ai également développé des compétences en programmation scientifique, notamment en Python, à travers plusieurs projets de travaux dirigés consacrés à l’analyse de données expérimentales et à la visualisation de résultats sous forme de courbes et de cartes de contours. Ces projets m’ont permis de développer une approche rigoureuse de l’analyse numérique et de la modélisation de phénomènes physiques. Par ailleurs, la rédaction de rapports scientifiques en LaTeX fait partie intégrante de ces travaux, renforçant mes compétences en rédaction scientifique.
			
			En parallèle de mes études, je suis membre de l’association TLS'e Racing, basée à l’ISAE-SUPAERO, qui développe un véhicule électrique pour la compétition Formula Student 2026. Dans ce cadre, j’ai participé au choix du moteur ainsi qu’à la définition de la configuration de la batterie. J’ai également réalisé des simulations structurelles avec Ansys afin d’évaluer la résistance du conteneur de batterie aux charges mécaniques et aux accélérations dans différentes directions, en respectant les contraintes du cahier des charges de la compétition.\\
						
			Je souhaiterais aujourd’hui poursuivre mon apprentissage dans un environnement stimulant où je pourrais mobiliser mes compétences en simulation numérique, renforcer ma compréhension de la mécanique des structures et contribuer activement à la réussite de projets techniques ambitieux.\\
			
			Veuillez agréer, Madame, Monsieur, l’expression de mes salutations distinguées.
	
	
	
	\vspace{0.8cm}
	
	
	
	
	
	
	\begin{flushright}
		
		\textbf{Friedly WOLI}
	\end{flushright}
	
	
\end{document}



























